% Esquisse d'un programme (cote 119, copy 2): archive cote 119, site page 107 (PDF page 108), typescript page 29.
% typescript. Draft compared: esquisse/pages/esq_119_031.tex (copy 1, site 31).
% AI-transcribed by Claude Opus 5.5 (claude-opus-5-5, adaptive thinking, effort high), Message Batches API, batch esq119_full (msgbatch_015xqDHkTumDysRF1qngZhRJ), 2026-11-13; not yet checked by eye.
\scanpage{119/107}
\mspage{29}

saires au développement d'un contexte utilisable de ``topologie
modérée''. Du coup cela relance une réflexion sur les fondements
d'une telle topologie, dont le besoin m'apparaît de plus en plus
clairement.

Avec un recul d'une dizaine d'années, je dirais aujour-%
d'hui, à ce sujet, que la ``\uline{topologie générale}'' \uline{a été développée}
(dans les années trente et quarante) \uline{par des analystes et pour}
\uline{les besoins de l'analyse}, non pour les besoins de la topologie
proprement dite, c'est à dire l'étude des \uline{propriétés topologiques}
\uline{de formes géométriques} diverses. Ce caractère inadéquat des fon-%
dements de la topologie se manifeste dès les débuts, par des
``faux problèmes'' (au point de vue au moins de l'intuition topo-%
logique des formes) comme celle de ``l'invariance du domaine'',
alors même que la solution de ce dernier par Brouwer l'amène à
introduire des idées géométriques nouvelles importantes. Aujour-%
d'hui encore, comme aux temps héroïques où on voyait pour la pre-%
mière fois et avec inquiétude des courbes remplir allègrement
des carrés et des cubes, quand on se propose de faire de la géo-%
métrie topologique dans le contexte technique des espaces topo-%
logiques, on se heurte à chaque pas à des difficultés parasites
tenant aux phénomènes sauvages. Ainsi, en dehors de cas de (très)
basse dimension, il ne peut guère être possible, pour un espace
donné $X$ (une variété compacte disons), d'étudier le type d'homo-%
topie (disons) du groupe des automorphismes de $X$, ou de l'espace
des plongements, ou immersions etc. de $X$ dans quelque autre
espace $Y$ - alors qu'on sent que ces invariants devraient faire
partie de l'arsenal des invariants essentiels associés à $X$, ou
au couple $(X,Y)$, etc, au même titre que l'espace fonctionnel
$\underline{\mathrm{Hom}}\,(X,Y)$ familier en topologie homotopique. Les topologues élu-%
dent la difficulté, sans l'affronter, en se rabattant sur des con-%
textes voisins du contexte topologique et moins marqués de sauva-%
gerie que lui, comme les variétés différentiables, les espaces PL
(linéaires par morceaux), etc, dont visiblement aucun n'est ``bon'',
i.e. n'est stable par les opérations topologiques les plus évidentes,

\margin{107}
